\section{Hydrogen transition and SADAR-lock atlas}
\label{sec:manual2-hydrogen-transition-sadar-lock-atlas}

This Hydrogen lane is the first Manual-II authoring pass after the H-1 native
occurrence gate.  Manual I materialized a target-blind H-1 direct-return
occurrence.  Manual II now carries that occurrence into a transition atlas by
freezing native packets first and keeping every metric, target, observation
join, notation link, and projection link outside the native AFC row; no target value is read before freeze.

\aodnoteblock{Authoring boundary}{The native packet is a finite ledger object:
identity, boundary, window, event order, returned-current relation, reclosure
relation, and SADAR-lock declaration.  An SI report can reference the packet only
after this freeze through a downstream index.  The report's temporal unit is the
second, written as \(\mathrm{s}\), and it is not a native premise.}

\subsection{Native packet order}

The deterministic generator is
\texttt{manual-2/scripts/build\_hydrogen\_transition\_sadar\_lock\_atlas.py}.
It binds the carried H-1 detector, trace, and occurrence rows and emits the
native transition packet ledger
\texttt{manual-2/data/hydrogen\_transition/hydrogen\_transition\_native\_packets.csv}.
The admitted native order is
\[
\text{H-1 returned-current anchor}
\to
RD
\to
RCD
\to
\SADAR\text{-lock declaration packet}.
\]
The optional SI-second projection contract is not part of this native order and
is not a native hash input.  The H-1 trace still reads
\[
\mathtt{H1\_q4\_state\_origin}
>
\mathtt{H1\_q4\_state\_hinge}
>
\mathtt{H1\_q4\_state\_origin}.
\]
No Balmer line, empirical wavelength, PubChem record, RDKit graph, UniProt row,
PDB coordinate, residual, score, Tau display symbol, projection identifier, or
SI scalar is an input to the native packet.


\subsection{Native RD/RCD, pressure, SADAR flow, and phase lock packets}

The Hydrogen transition native-materialization pass extends the prior freeze
ledger with four exact native packets.  These packets are generated by the
Hydrogen native-materialization script and remain target-blind; the concrete
script path is recorded in the source manifest.

\begin{center}
\begin{tabularx}{\linewidth}{@{}lX@{}}
\toprule
Ledger & Role \\
\midrule
\texttt{\detokenize{hydrogen_transition_rd_rcd_packets.csv}} & exact single-path RD distribution and RCD coupling for the carried H-1 direct-return trace. \\
\texttt{\detokenize{hydrogen_transition_duonic_pressure_packets.csv}} & local duonic-pressure packet; pressure is a coupling load, not cadence. \\
\texttt{\detokenize{hydrogen_transition_sadar_flow_packets.csv}} & native SADAR-flow packet; SADAR flow is not a metric-time report. \\
\texttt{\detokenize{hydrogen_transition_phase_lock_packets.csv}} & subject/reference phase-lock packet with metric reference still undeclared. \\
\bottomrule
\end{tabularx}
\end{center}

The materialized values are exact rational rows.  The direct returned-current
trace has one admitted path of mass one, with two executed bip tokens recorded as
execution structure rather than temporal magnitude.  The RD packet records
\[
RD=2,
\]
for this one path, and the RCD packet couples that returned duration to the
outbound/return asymmetry signature
\[
+1|-1.
\]
The duonic-pressure packet uses
\[
\rho^D_\omega=2,
\qquad
C_{\mathrm{local}}=1,
\qquad
p^D=2.
\]
The SADAR-flow packet records
\[
A=1,
\qquad
\mathfrak S=2,
\]
with status
\texttt{native\_SADAR\_flow\_not\_metric\_time}.  The phase-lock packet records
a native self-reference lock
\[
1:1,
\qquad
\Delta_\phi=0,
\]
while the metrological reference remains undeclared and the metric-time report
remains unmaterialized.

\aodnoteblock{Native-flow boundary}{This materialization separates flow from
measurement.  RD/RCD, duonic pressure, and SADAR flow are native current packets.
A temporal measurement requires a subject/reference relational lock.  A metric
second requires a later metrological reference card.}

The counterfactual audit for this pass rejects target insertion into RD/RCD,
pressure-as-cadence relabeling, premature metric-time materialization, and early
metrological-reference declaration.  The unchanged control passes.

\subsection{Hash-isolated projection and notation indices}

The native packet hash is computed only from native columns: packet order,
packet identity, source binding, H-1 occurrence and detector identifiers,
boundary, window, native event order, freeze status, target-read status,
observable-join state, and empirical-score status.  Projection and notation
metadata live in separate ledgers:
\begin{center}
\begin{tabularx}{\linewidth}{@{}lX@{}}
\toprule
Ledger & Role \\
\midrule
\texttt{\detokenize{hydrogen_transition_projection_index.csv}} & links a frozen native
packet to an optional downstream SI-second projection packet; changing the
projection identifier changes only the projection-index row hash. \\
\texttt{\detokenize{hydrogen_transition_tau_notation_index.csv}} & links a frozen native
packet to a Tau display symbol; changing the display symbol changes only the
notation-index row hash. \\
\bottomrule
\end{tabularx}
\end{center}
Thus native packet identity is stable under optional projection or rendering
changes.

\subsection{Tau cycle notation}

Cycle reporting in this lane uses Tau notation.  The full scoped cycle is
recorded as
\[
\Tau_{H1}
:
\mathtt{origin} \to \mathtt{hinge} \to \mathtt{origin},
\qquad
\Tau_{\mathrm{lock}}
:
\text{freeze before optional projection}.
\]
This notation is a ledger convention for full-cycle reporting.  It does not add
a circle constant to the AFC primitive layer, and it does not convert the two
executed bip tokens into a metric duration.  The registry file
\texttt{tau\_cycle\_notation\_registry.csv} records this as a Tau reporting rule
and marks the circle-constant alternative as forbidden as a native primitive.

\subsection{SI-second projection lane}

The metric anchoring lane is intentionally downstream.  Its SI-second
projection-packet ledger lives in the Hydrogen-transition data directory.  The
active rows say:
\begin{center}
\begin{tabular}{@{}ll@{}}
projection layer & \texttt{downstream\_metric\_report\_only} \\
SI temporal unit & \texttt{second} \\
map state & \texttt{contract\_declared\_not\_instantiated}
\end{tabular}
\end{center}
Thus the native row may be audited without a clock.  A later SI report may use
\(\mathrm{s}\) as the temporal unit, but the native Hydrogen transition packet
continues to carry only event order, boundary, window, and native packet hashes.

\subsection{Matter-octave SADAR-lock atlas}

The atlas also names the Manual-II matter scales that inherit the same freeze
discipline:
\begin{center}
\begin{tabularx}{\linewidth}{@{}lXl@{}}
\toprule
Lane & Bound source & Status \\
\midrule
Elementary Matter 118 & \texttt{element\_registry\_118.csv} plus the
\(3{:}3{:}6\) and \(3{:}4{:}6\) ladders & typed atlas row hash-locked \\
Molecular Matter & \texttt{formula\_residuals.csv} and chain formula packets &
typed atlas row hash-locked \\
Biomolecular Matter & nucleotide, peptide, and protein-chain candidate rows &
typed atlas row hash-locked \\
\bottomrule
\end{tabularx}
\end{center}
The atlas file is
\texttt{sadar\_lock\_matter\_octave\_atlas.csv}.  Its elementary row binds all
118 element identity rows.  Its molecular row binds exact CHONPS residual
closures.  Its biomolecular row binds nucleotide, peptide, and short protein
chain closures.  External lookup lanes remain named but closed: PubChem, RDKit,
UniProt, PDB, AlphaFold, and SI-second reports do not enter the native atlas
rows.

\subsection{Executed fail-closed audit}

The counterfactual audit is executed by the generator before the CSV is written.
The evaluator mutates native target insertion, SI-second promotion into the
native row, circle-constant promotion into the native cycle notation, and
external target-lane join before packet freeze.  Each mutated case is rejected.
The unchanged control passes.  Therefore the authoring claim is narrow.
\aodnoteblock{Authoring claim}{Hydrogen materializes target-blind native H-1
transition declaration packets and matter-octave atlas rows.  SI-second
reporting, empirical targets, residuals, and scores remain downstream and
closed.}
Independent review is still required before this authoring result can advance
to GO.
